(1) 膜的小挠度满足 \(\rho_s u_{tt}=T_xu_{xx}+T_yu_{yy}\)，边界 \(u(0,y)=u(L,y)=u(x,0)=u(x,L)=0\)。各向异性张力方向固定，故仍可用变量分离，且两方向导数权重不同。
(2) \(u_{mn}=\sin(m\pi x/L)\sin(n\pi y/L)e^{i\omega_{mn}t}\)，\(m,n\ge1\)，代入得 \(\omega_{mn}^2=\pi^2[T_xm^2+T_yn^2]/(\rho_sL^2)\)，每一模的动能和张力势能均为二次型。
(3) 记 \(C=\pi^2T/(\rho_sL^2)\)，则 \(\omega_{12}^2=C(5-3\epsilon)\)、\(\omega_{21}^2=C(5+3\epsilon)\)，所以 \(\omega_{21}^2-\omega_{12}^2=6C\epsilon\)。在 \(\epsilon=0\) 两模简并，张力各向异性破坏交换 \(x\leftrightarrow y\) 对称性。
(4) 总能量 \(\mathcal E=\frac12\int[\rho_su_t^2+T_xu_x^2+T_yu_y^2]\,dA\)，当 \(|\epsilon|<1\) 时严格正定且 \(d\mathcal E/dt=0\)。\(\epsilon\to0\) 恢复各向同性简并；\(|\epsilon|\to1\) 时某方向张力趋零，模型的二维膜稳定性变为退化情形。
